\documentclass[11pt]{article}
\usepackage[margin=1in]{geometry}
\usepackage{amsmath, amssymb, amsthm}
\usepackage{hyperref}

\newcommand{\blank}{\underline{\hspace{1.5cm}}}

\title{MAT 3150 --- Homework 3\\
\large Problem 4 Walkthrough: Averages of an Increasing Sequence}
\author{}
\date{}

\begin{document}
\maketitle

\textbf{Goal.} Let $(x_n)_n$ be an increasing sequence of positive real
numbers and define
\[
  y_n=\frac{x_1+x_2+\cdots+x_n}{n}.
\]
Prove that $(y_n)$ is increasing as well.

% ============================================================
\section*{What ``increasing'' means}
A sequence $(a_n)$ is increasing if $a_n \le a_{n+1}$ for every $n$.
(If your course uses \emph{strictly} increasing, replace $\le$ by $<$
everywhere below; the argument is the same.)

\textbf{The standard way to prove it:} show that
\[
  a_{n+1} - a_n \ge 0 \quad\text{for every } n .
\]
With fractions, that means: put $a_{n+1}-a_n$ over a common denominator,
note the denominator is positive, and show the numerator is $\ge 0$.

\textbf{Notation that makes this problem manageable.} Let
\[
  S_n = x_1 + x_2 + \cdots + x_n, \qquad\text{so}\qquad y_n = \frac{S_n}{n}.
\]
The one fact about $S_n$ you'll need:
\[
  S_{n+1} = S_n + x_{n+1}.
\]

% ============================================================
\section*{Worked example: $a_n = \dfrac{n}{n+1}$ is increasing}

\paragraph{Compute the difference.} Common denominator $(n+1)(n+2)$:
\[
  a_{n+1}-a_n = \frac{n+1}{n+2}-\frac{n}{n+1}
  = \frac{(n+1)^2 - n(n+2)}{(n+1)(n+2)}
  = \frac{n^2+2n+1 - n^2-2n}{(n+1)(n+2)}
  = \frac{1}{(n+1)(n+2)} .
\]

\paragraph{Sign check.} The numerator $1>0$ and the denominator
$(n+1)(n+2)>0$, so $a_{n+1}-a_n>0$.

\begin{proof}
For every $n\in\mathbb{N}$,
\[
  a_{n+1}-a_n = \frac{1}{(n+1)(n+2)} > 0,
\]
since the numerator and denominator are both positive. Hence
$a_n < a_{n+1}$ for every $n$, so $(a_n)$ is increasing.
\end{proof}

\noindent\textbf{Pattern:} difference $\to$ common denominator $\to$
simplify $\to$ show numerator and denominator have the right signs.

% ============================================================
\newpage
\section*{Your turn: $y_n = S_n/n$}

\paragraph{1. Set up the difference. \checkmark\ (done)}
\[
  y_{n+1}-y_n = \frac{S_{n+1}}{n+1} - \frac{S_n}{n} .
\]

\paragraph{2. Common denominator. \checkmark\ (done)} Multiply the first
fraction by $\frac{n}{n}$ and the second by $\frac{n+1}{n+1}$:
\[
  y_{n+1}-y_n = \frac{n\,S_{n+1} - (n+1)\,S_n}{n(n+1)} .
\]

\paragraph{Denominator sign. \checkmark\ (done)} $n(n+1)>0$ since
$n\ge 1$. So the sign of $y_{n+1}-y_n$ is the sign of the numerator.

\paragraph{Expand the numerator. \checkmark\ (done)} Substitute
$S_{n+1} = S_n + x_{n+1}$:
\[
  n\,S_{n+1} - (n+1)\,S_n
  = \bigl(n\,S_n + n\,x_{n+1}\bigr) - \bigl(n\,S_n + S_n\bigr).
\]

\paragraph{3. Cancel. \checkmark\ (done)} Distribute the minus
sign over both terms in the second bracket:
\[
  \bigl(n\,S_n + n\,x_{n+1}\bigr) - \bigl(n\,S_n + S_n\bigr)
  = n\,S_n + n\,x_{n+1} - n\,S_n - S_n .
\]
The first term $+n\,S_n$ and the third term $-n\,S_n$ cancel
\checkmark. The two terms that remain:
\[
  = n\,x_{n+1} - S_n .
\]
So
\[
  y_{n+1}-y_n = \frac{n\,x_{n+1} - S_n}{n(n+1)} .
\]

\bigskip
\hrule
\bigskip
\noindent\textbf{\large Step 4 --- where ``increasing'' is used.}

\paragraph{Numerical example first.} Take the increasing sequence
$x_1=2,\ x_2=5,\ x_3=7,\ x_4=10$ and $n=3$, so $x_{n+1}=x_4=10$.
\begin{itemize}
  \item Each earlier term is at most the new one:
        $2\le 10,\quad 5\le 10,\quad 7\le 10$.
  \item Add them:
        $\underbrace{2+5+7}_{S_3=14} \le \underbrace{10+10+10}_{3\cdot 10=30}$.
  \item Numerator: $3\cdot 10 - 14 = 16 \ge 0$.
  \item Check: $y_3 = 14/3 \approx 4.67$, $y_4 = 24/4 = 6$, and
        $\dfrac{16}{3\cdot 4} = \dfrac43 = y_4 - y_3$. \checkmark
\end{itemize}

\paragraph{Example $\to$ general.}
\begin{center}
\renewcommand{\arraystretch}{1.4}
\begin{tabular}{l|l|l}
  & Example ($n=3$) & General \\ \hline
  New term & $x_4 = 10$ & $x_{n+1}$ \\
  Earlier terms & $2,5,7$ & $x_1,\ldots,x_n$ \\
  Each earlier $\le$ new & $2\le10,\ 5\le10,\ 7\le10$
    & $x_k\le x_{n+1}$, $k=1,\ldots,n$ \\
  Sum of left sides & $2+5+7=S_3$ & $x_1+\cdots+x_n=S_n$ \\
  Sum of right sides & $10+10+10=3\cdot10$
    & $x_{n+1}+\cdots+x_{n+1}=n\,x_{n+1}$ \\
  Result & $S_3\le 3\cdot 10$ & $S_n\le n\,x_{n+1}$
\end{tabular}
\end{center}

\paragraph{4. Show the numerator is $\ge 0$. \checkmark\ (done)}
Since $(x_n)$ is increasing,
\[
  x_k \le x_{n+1} \quad\text{for each } k=1,\ldots,n.
\]
Adding these $n$ inequalities: the left sides sum to $S_n$, and the
right side is $n$ copies of $x_{n+1}$, so
\[
  S_n \le n\,x_{n+1} .
\]
Subtracting $S_n$ from both sides,
\[
  n\,x_{n+1} - S_n \ge 0 .
\]

\paragraph{5. Conclude. \checkmark\ (done)} Numerator $\ge 0$ and
denominator $n(n+1)>0$, so
\[
  y_{n+1}-y_n = \frac{n\,x_{n+1} - S_n}{n(n+1)} \ge 0,
\]
i.e.\ $y_n \le y_{n+1}$ for every $n$.

\paragraph{In your own words:} why is $S_n \le n\,x_{n+1}$? Finish the
sentence: ``Each of the $n$ terms $x_1,\ldots,x_n$ is at most \blank,
so their sum is at most \blank.''

\paragraph{6. The proof.} Fill in the template:
\begin{quote}
Let $S_n = x_1+\cdots+x_n$, so $y_n = S_n/n$ and
$S_{n+1} = S_n + x_{n+1}$. For every $n$,
\[
  y_{n+1}-y_n = \frac{\blank}{n(n+1)} .
\]
Since $(x_n)$ is increasing, $x_k \le x_{n+1}$ for $k=1,\ldots,n$;
adding these gives $S_n \le \blank$. Hence the numerator is $\ge 0$.
The denominator $n(n+1)$ is positive. Therefore
$y_{n+1}-y_n \ge 0$, i.e.\ $y_n \le y_{n+1}$ for every $n$, so $(y_n)$
is increasing.
\end{quote}

\paragraph{Check yourself.}
\begin{itemize}
  \item Where did you use that $(x_n)$ is increasing?
  \item Did you use that the $x_n$ are \emph{positive}? (It's fine if
        you didn't --- just notice.)
\end{itemize}

\end{document}
